\section{Related Work}\label{related-work}

The benchmark sits at the intersection of seven distinct literatures: static geometry of LLM representations, the theoretical origin of that geometry, conceptual navigation and interpolation, consistency and directionality, multi-hop reasoning, cross-model comparison, and cognitive science foundations. Each has produced results that inform our design and findings, but none asks the question we ask: whether models can \emph{navigate} conceptual space consistently and whether the resulting paths have measurable geometric structure. We review each literature and state precisely what gap the benchmark fills.

\subsection{Static Geometry of LLM Representations}\label{static-geometry-of-llm-representations}

The linear representation hypothesis \citep{park2024linear} established that individual concepts correspond to linear directions in activation space. Its extension \citep{park2025geometry} showed that hierarchical relationships form polytope structures --- convex regions whose geometric relationships encode taxonomic nesting. Separately, hyperbolic embedding methods \citep{nickel2017poincare} demonstrated that tree-like hierarchies are better captured in negatively curved spaces, a finding extended to LLM representations by recent work on HELM~\citep{helm2025} and direct measurement of delta-hyperbolicity and ultrametricity~\citep{deltahyperbolicity2025}, which quantify how tree-like a given embedding space actually is.

Topological data analysis has added finer-grained structural characterization. Explainable Mapper~\citep{explainablemapper2025} identified circular manifolds for temporal concepts and linear manifolds for ordered quantities in LLM representations. Persistent Topological Features~\citep{persistenttopological2025} used persistence diagrams to characterize the homological structure of concept clusters. \citet{wyss2025spectral} proved mathematically that LLM activation manifolds partition into discrete basins of attraction via spectral analysis of semantic attractors --- a result directly relevant to the bottleneck bridge structure we observe empirically (Section 5.4). \citet{joshi2025dimensionality} characterized intrinsic dimensionality profiles across transformer layers, finding a universal early-expand-late-compress pattern that holds across architectures.

\textbf{Gap we fill.} This entire literature characterizes \emph{what shape the space has} --- its curvature, dimensionality, manifold topology, and basin structure. All of it is static. We probe whether models can \emph{navigate} that space consistently, producing paths with measurable geometric properties. The distinction is between a map and a route: static geometry gives us the map; our benchmark tests whether models can walk it.

\subsection{Theoretical Origin of Geometric Structure}\label{theoretical-origin-of-geometric-structure}

\citet{karkada2025translation} proved that translation symmetry in word co-occurrence statistics mathematically determines manifold geometry in embedding spaces. Their result is architecture-independent: the same geometric structures arise in Word2Vec, GloVe, and Gemma because all train on data distributions with shared statistical regularities. Geometry comes from data, not architecture.

This theoretical result connects directly to our empirical findings. If geometric structure is determined by data statistics that all models share, then navigational structure derived from that geometry should also generalize across architectures --- which is what we observe (R1 and R2 universal across 12 models from 11 training pipelines; Section 8). The scale effect (Llama 8B as outlier in bridge frequency) is consistent with their framework: if extracting the full geometric structure requires sufficient capacity to represent translation symmetries, smaller models may have the coarse structure but lack the resolution to navigate through the same landmarks as frontier models.

\textbf{Gap we fill.} Karkada et al.~study the static manifold shape and its mathematical origin. We test the dynamic complement: whether that geometry supports consistent behavioral navigation. Their theory is consistent with our empirical finding of cross-architecture generalization; our data provides the behavioral evidence their framework implies but does not test. The contributions are complementary --- they explain why the geometry exists; we show what it supports.

\subsection{Conceptual Navigation and Interpolation}\label{conceptual-navigation-and-interpolation}

The mathematical framework for navigation through latent spaces exists but has not been applied to LLM behavioral outputs. \citet{arvanitidis2018latent} developed Riemannian geodesics as semantically meaningful paths in generative model latent spaces --- shortest paths through curved space that respect the local geometry. The framework is directly applicable to conceptual navigation but has been used for image generation and VAE interpolation, not for probing how language models traverse between concepts.

The closest precedent for waypoint-based navigation is Relation Embedding Chains~\citep{relationchains2023}, which discovered intermediate concepts X such that A-to-B can be decomposed as A-to-X-to-B using knowledge graph paths. This is the nearest prior work to our waypoint elicitation paradigm, though it operates on static knowledge graph structure rather than behavioral generation. The Geometry of Knowledge~\citep{geometryknowledge2025} traversed semantic manifolds for generative diversity --- the closest existing work to ``conceptual traversal'' --- but used traversal as a tool for improving generation, not as an evaluation paradigm.

More recently, reasoning-as-path-search has entered the mainstream. RiemannInfer~\citep{riemanninfer2026} frames multi-step inference as trajectory optimization on a Riemannian manifold, and Geometric Reasoner~\citep{geometricreasoner2026} uses geodesic metaphors for chain-of-thought planning. These works treat the path metaphor as productive for reasoning but do not test the geometric properties of the paths themselves.

\textbf{Gap we fill.} No prior work asks models to produce intermediate waypoints between two concepts and then tests whether the resulting paths have consistent geometric structure --- satisfying metric axioms, exhibiting compositional properties, and responding to causal perturbation. The mathematical framework exists (Riemannian geodesics); the precedent for intermediate concepts exists (relation embedding chains); the path metaphor is becoming mainstream (reasoning-as-search). What is missing is the empirical test of whether behaviorally generated paths have measurable, replicable geometry. We provide that test.

\subsection{Consistency, Directionality, and the Reversal Curse}\label{consistency-directionality-and-the-reversal-curse}

\citet{berglund2024reversal} established the reversal curse: autoregressive models trained on ``A is B'' systematically fail to learn ``B is A.'' This is an architectural limitation --- the left-to-right generation order creates inherent directional bias. the directional optimization asymmetry literature~\citep{directionalasymmetry2025} deepened this finding by showing that inverse mappings incur higher excess loss even under symmetric training objectives, establishing that the A-to-B versus B-to-A gap is architectural, not merely a training-data artifact. \citet{barakat2026hysteresis} extended directionality to semantic caching, demonstrating that input ordering steers inference trajectories into different attractor basins --- directional hysteresis at the system level.

On the consistency side, recent work has documented pervasive self-consistency failures in frontier models~\citep{selfconsistency2025}, with high variability even in GPT-4 across repeated evaluations on identical prompts. Trajectory variance in agentic settings~\citep{trajectoryvariance2025} shows substantial action-sequence diversity across runs of the same agent on the same task, establishing that behavioral variability is not confined to simple generation tasks.

\textbf{Our contribution.} Existing work treats inconsistency and asymmetry as failure modes --- problems to be solved or architectural limitations to be documented. We reframe asymmetry as data. The pervasive directional asymmetry we observe (mean 0.811 in the Phase 2 cohort, exceeding 0.60 across all 12 models tested) is not a bug but a structural property of navigational space, consistent with quasimetric topology. The pattern of asymmetry reveals structure: it tracks semantic properties of concept pairs and model-specific navigational styles, producing systematic geometric signatures rather than noise.

A defensive clarification is warranted here. Reviewers will ask why this is not simply self-consistency work with a geometry gloss. The contrast is precise: self-consistency work measures the \emph{magnitude} of variance across samples and treats it as error to be minimized. We test whether the \emph{structure} of that variance satisfies formal geometric axioms --- metric properties, compositionality, causal sensitivity. We do not ask whether models are consistent; we ask whether the pattern of their inconsistency has geometric regularity. Variance is our data, not our failure mode.

\subsection{Multi-Hop Reasoning and Compositionality}\label{multi-hop-reasoning-and-compositionality}

The compositionality gap \citep{press2022compositionality} demonstrated that models can answer individual sub-questions correctly while failing the composed question --- knowing both ``A is B'' and ``B is C'' but failing ``A is C.'' the two-hop curse~\citep{twohopcurse2025} formalized this for knowledge composition: models trained on A-to-B and B-to-C facts fail A-to-C inference without explicit chain-of-thought prompting. TWOHOPFACT~\citep{twohopfact2024} studied latent multi-hop reasoning through bridge entities, identifying the intermediate step as the critical point of failure --- the model must internally activate the bridge concept to compose the two hops.

\textbf{Our contribution.} We generalize the compositionality question from binary pass/fail to geometric characterization. Multi-hop reasoning work asks whether the model gets the composed answer \emph{correct}. We ask what \emph{shape} the intermediate sequence takes --- whether waypoints compose consistently (hierarchical transitivity 4.9x over random controls), whether bridge concepts function as structural bottlenecks (frequencies of 0.60-1.00 for predicted bridges), and whether the compositional structure satisfies the triangle inequality (\textasciitilde91\% across three independent samples). Waypoints externalize latent composition, making the intermediate steps available for geometric analysis rather than accuracy scoring.

A second defensive clarification applies here. The paradigm will invite comparison to chain-of-thought elicitation. The contrast is fundamental: chain-of-thought elicits reasoning steps toward a \emph{correct answer} --- intermediate steps are evaluated by whether they lead to a right response. We elicit navigational waypoints toward a \emph{destination concept} --- there is no correct answer, only geometric structure to characterize. Chain-of-thought is evaluated by task accuracy; our paradigm is evaluated by structural consistency, metric properties, and causal sensitivity. The two paradigms share surface similarity (both produce ordered intermediate steps) but differ in what they measure and how they are validated.

\subsection{Cross-Model Comparison and the Platonic Representation Hypothesis}\label{cross-model-comparison-and-the-platonic-representation-hypothesis}

The Platonic Representation Hypothesis \citep{huh2024platonic} predicts that model representations converge with scale, measured by CKA alignment between activation spaces. Larger models from different families produce increasingly similar internal representations --- a convergence toward a shared ``platonic'' geometry. Multi-way alignment work~\citep{multiwayalignment2026} extended this to shared reference spaces across N models, and model stitching experiments~\citep{modelstitching2025} demonstrated that features transfer meaningfully but non-uniformly across architectures.

On the behavioral side, LLM behavioral fingerprinting~\citep{behavioralfingerprint2025} established that models have characteristic behavioral signatures --- stable patterns of response that distinguish one model from another across tasks. \citet{nanda2026attractors} identified conversational attractor states: model-specific behavioral loops (Claude trending toward philosophical synthesis, GPT toward system-building) that function as the conversational analog of our quantitative gaits. Behavioral failure manifold mapping~\citep{failuremanifold2025} characterized model-specific topographies with qualitatively different basin structures, showing that models fail in geometrically distinct ways.

\textbf{Our contribution.} The PRH measures static representational alignment. We test \emph{dynamic navigational alignment} --- whether shared representations imply shared routes. The finding is a three-level hierarchy: structure generalizes (all 12 models have gait and asymmetry), content mostly generalizes among frontier-scale models (aggregate bridge frequency difference CI includes zero), but gaits diverge (the 2.9x gait range does not collapse with scale). Same map, different routes. This is a behavioral refinement of the PRH: representational convergence does not imply behavioral convergence. The models share structural properties (gait and asymmetry are universal) and largely converge on the strongest navigational landmarks, but they navigate with model-specific route preferences that remain distinctive even at frontier scale --- and their navigational distances are fundamentally model-dependent (Section 7.4).

\subsection{Cognitive Science Foundations}\label{cognitive-science-foundations}

\citet{gardenfors2000} proposed Conceptual Spaces as geometric similarity spaces organized along quality dimensions, with natural concepts forming convex regions and betweenness relations defining navigable paths. This theoretical framework posits that conceptual structure is inherently geometric --- concepts have locations, distances between them are meaningful, and navigation between concepts follows principled geometric trajectories.

Cognitive neuroscience has found biological mechanisms consistent with this theory. \citet{tolman1948cognitive} proposed cognitive maps for spatial navigation. \citet{okeefe1971hippocampus} discovered place cells in the hippocampus. \citet{constantinescu2016grid} demonstrated that grid cells in the human entorhinal cortex fire during \emph{conceptual} navigation --- not just physical movement through space but traversal through abstract conceptual dimensions. The spatial navigation machinery is repurposed for abstract thought.

Two additional cognitive science results bear directly on our findings. \citet{tversky1977features} established that human similarity judgments are asymmetric: North Korea is judged more similar to China than China is to North Korea. This is the direct cognitive parallel to our quasimetric finding --- asymmetry in conceptual distance is not an artifact of autoregressive generation but a property also observed in human cognition. \citet{collins1975spreading} proposed spreading activation through semantic networks, where activation propagates directionally from a source concept through associative links --- a process model that maps onto the early-anchoring and bridge-bottleneck phenomena we observe.

More recently, LLM World of Words~\citep{llmwow2025} demonstrated that free association norms elicited from LLMs are directly comparable to human norms collected over decades, establishing that LLM conceptual organization shares at least surface-level structure with human semantic memory.

\textbf{Our contribution.} Gardenfors theorized that concepts live in navigable geometric spaces. Cognitive neuroscience found biological mechanisms (grid cells, place cells) that support navigation through those spaces. We provide a computational probe of this theoretical framework at scale --- 12 models, 100+ concept pairs, \textasciitilde21,540 runs. The partial confirmation is substantive: LLM behavior is consistent with concepts inhabiting navigable spaces with measurable, compositional distance structure (triangle inequality \textasciitilde91\%, hierarchical transitivity 4.9x over random). But the pervasive asymmetry (mean 0.811) challenges the standard metric assumption in conceptual spaces theory, suggesting that quasimetric structure may be a more appropriate formal framework. We do not claim that LLMs replicate human conceptual navigation --- only that they exhibit analogous geometric regularities testable by similar methods. The parallel to Tversky's asymmetry and Constantinescu's grid-cell navigation is structural, not mechanistic.

\begin{center}\rule{0.5\linewidth}{0.5pt}\end{center}

The seven literatures reviewed here converge on a single gap. Static geometry is well-characterized and its theoretical origin is understood. The mathematical framework for navigation exists. Consistency and directionality are documented as phenomena. Multi-hop composition is studied as a reasoning failure mode. Cross-model convergence is measured at the representational level. Cognitive science provides both the theoretical framework and the biological precedent. What is missing --- and what this benchmark provides --- is an empirical test of whether LLM conceptual navigation has consistent, measurable geometric structure: not what the space looks like, but whether models can walk through it reliably, and what the shape of their walking reveals.
